% !TEX root = ../../../main.tex
% 来源：中学数学实验教材·第二册（几何）
% 章节：圆 / 圆的基本性质
% 原始文件：4.tex，行 894-947
% 类型：tikz
% ---
  \begin{tikzpicture}[scale=1.3]
  \begin{scope}
    \tkzDefPoints{0/0/O_1, 1.6/0/O_2, .5/0/a, .8/0/b}
  \tkzDrawCircles[thick](O_1,a O_2,b)
  \tkzDrawSegments[thick](O_1,O_2)
    \tkzDrawPoints(O_1,O_2)
    \tkzLabelPoints[below](O_1,O_2)
    \node at (1,-1){(1)};
  \end{scope}
  \begin{scope}[xshift=3.5cm]
    \tkzDefPoints{0/0/O_1, 1.3/0/O_2, .5/0/a}
  \tkzDrawCircles[thick](O_1,a O_2,a)
  \tkzDrawSegments[thick](O_1,O_2)
    \tkzDrawPoints(O_1,O_2)
    \node at (.7,-1){(2)};
    \tkzLabelPoints[below](O_1,O_2)
  \end{scope}
  \begin{scope}[xshift=7cm]
    \tkzDefPoints{0/0/O_1, 1.1/0/O_2, .5/0/a, .3/0/b}
    \tkzDrawCircles[thick](O_1,a O_2,b)
    \tkzDrawSegments[thick](O_1,O_2)
    \tkzLabelPoints[below](O_1,O_2)
    \tkzDrawPoints(O_1,O_2)
  \tkzInterCC(O_1,a)(O_2,b) \tkzGetPoints{A}{B}
  \tkzLabelPoints[above](A)
  \tkzLabelPoints[below](B)
    \node at (.5,-1){(3)};
    \tkzDrawSegments[dashed](O_1,A O_2,A)
  \tkzDrawSegments[thick](A,B)
  \end{scope}
  \begin{scope}[yshift=-2.5cm, xshift=1cm]
    \tkzDefPoints{0/0/O_1, .6/0/O_2, -.5/0/a}
    \tkzDrawCircles[thick](O_1,a O_2,a)
    \tkzDrawSegments[thick](a,O_2)
      \tkzDrawPoints(O_1,O_2)
      \tkzLabelPoints[below](O_1,O_2)
      \node at (.6,-1.5){(4)};
  \end{scope}
  \begin{scope}[yshift=-2.5cm, xshift=4cm]
    \tkzDefPoints{0.2/0/O_1, .6/0/O_2, -.5/0/a, -.4/0/b}
    \tkzDrawCircles[thick](O_1,b O_2,a)
    \tkzDrawSegments[thick](O_1,O_2)
      \tkzDrawPoints(O_1,O_2)
      \tkzLabelPoints[below](O_1,O_2)
      \node at (.6,-1.5){(5)};
  \end{scope}
  \begin{scope}[yshift=-2.5cm, xshift=7.5cm]
    \tkzDefPoints{0/0/O, 1/0/a, .7/0/b}
    \tkzDrawCircles[thick](O,b O,a)
      \tkzDrawPoints(O)
      \tkzLabelPoint[below](O){$O_1(O_2)$}
      \node at (0,-1.5){(6)};
  \end{scope}
  \end{tikzpicture}
